\documentclass[11pt,a4paper]{article}
\usepackage[utf8]{inputenc}
\usepackage[T1]{fontenc}
\usepackage[margin=0.75in]{geometry}
\usepackage{enumitem}
\usepackage{titlesec}
\usepackage{xcolor}
\usepackage{hyperref}

\definecolor{primarycolor}{RGB}{33,37,41}
\definecolor{secondarycolor}{RGB}{0,102,204}
\definecolor{accentcolor}{RGB}{108,117,125}
\hypersetup{colorlinks=true,urlcolor=secondarycolor}
\titleformat{\section}{\Large\bfseries\color{primarycolor}}{}{0em}{}
\titlespacing*{\section}{0pt}{1.2em}{0.5em}
\setlist[itemize]{leftmargin=*,itemsep=0.2em,topsep=0.25em}
\pagestyle{empty}

\begin{document}

\begin{center}
{\Huge\bfseries\color{primarycolor} MARTA KOWAL-BRANDT}\\[0.35em]
{\large\color{accentcolor} QA Automation Engineer | CI/CD and Test Infrastructure}\\[0.25em]
Wroclaw, Poland | marta.kowalbrandt@example.com | +48 555 303 303 | gitlab.com/mkb-example
\end{center}

\section*{PROFESSIONAL SUMMARY}
QA automation engineer with 7 years of experience, including recent hands-on work building and maintaining CI/CD pipelines, test environments, and build tooling. Experienced with GitLab CI, Docker, Helm, Kubernetes-based test environments, Linux, and Python and Bash automation. Moving toward DevOps and platform engineering, with additional personal experience running a k3s homelab and provisioning a Hetzner VPS with Terraform.

\section*{TECHNICAL SKILLS}
CI/CD and Infrastructure: GitLab CI, Jenkins, Docker, Docker Compose, Helm, Terraform (personal projects)\\
Containers and Platforms: Kubernetes (Helm deployments and namespace-level access), k3s, Argo CD\\
Monitoring and Data: Prometheus, Grafana, Loki, PostgreSQL\\
Programming and Scripting: Python, Bash, Java (test code), SQL\\
Testing: pytest, Playwright, Selenium, Postman\\
Systems: Linux, Git, Ansible (basic)\\
Cloud learning: AWS EC2 and S3 basics

\section*{PROFESSIONAL EXPERIENCE}
\textbf{QA Automation Engineer (Test Infrastructure) | Odrava Health Software | Wroclaw} \hfill 2021-Present
\begin{itemize}
\item Own GitLab CI pipelines for 3 product teams (approximately 25 developers), covering build, unit, integration, and end-to-end test stages and artifact publishing.
\item Reduced average pipeline time from 38 to 14 minutes through parallel test jobs, dependency caching, and separating end-to-end tests into their own stage; reduced monthly CI runner costs by about 30\%.
\item Built short-lived merge-request test environments with Docker Compose and later Helm on the shared Kubernetes cluster; environments include the application, PostgreSQL, and seeded data, and are automatically removed after merge or 48 hours of inactivity.
\item Maintain a pytest and Playwright end-to-end suite of approximately 900 tests; reduced the flaky-test rate from 7\% to below 1\%.
\item Wrote Bash and Python tooling to seed PostgreSQL test data and built Grafana dashboards using CI test results to help teams identify flaky tests.
\item Wrote the internal guide to how CI works, used in onboarding; mentored 2 QA interns in 2023.
\item Work with namespace-level access on the shared Kubernetes cluster; the platform team operates the cluster.
\end{itemize}

\textbf{QA Engineer | Brightmill Software House | Wroclaw} \hfill 2019-2021
\begin{itemize}
\item Tested web and mobile applications for client projects using manual and automated testing.
\item Introduced a Java and Selenium automated regression suite for the largest client project.
\item Configured Jenkins jobs for nightly regression runs.
\end{itemize}

\textbf{Customer Support Specialist (2nd line) | Brightmill Software House | Wroclaw} \hfill 2018-2019
\begin{itemize}
\item Investigated customer issues using logs and databases, and reproduced bugs for developers.
\end{itemize}

\textbf{Homelab k3s Cluster with GitOps | Personal} \hfill 2024-2025
\begin{itemize}
\item Ran a three-node k3s cluster on Raspberry Pi 5 devices; deployed applications from a Git repository using Argo CD.
\item Set up Prometheus, Grafana, and Loki for monitoring, with Traefik ingress and cert-manager for TLS.
\item Provisioned infrastructure for a small Hetzner VPS with Terraform; used Ansible at a basic level.
\item Technologies: k3s, Argo CD, Prometheus, Grafana, Loki, Terraform, Ansible.
\end{itemize}

\section*{EDUCATION}
BEng, Computer Engineering | Silesian Lowlands University of Technology | 2018

\section*{CERTIFICATIONS AND DEVELOPMENT}
ISTQB Foundation Level | 2019\\
Certified Kubernetes Administrator (CKA) preparation | Exam scheduled for January 2027

\section*{LANGUAGES}
Polish (native), English (C1), German (A2)

\end{document}